\addchap{Abstract}

Entity typing assigns semantic categories to entities. When categories form a taxonomy, a useful prediction must identify a suitable specific type, recover relevant broader types, and distinguish nearby alternatives. This thesis investigates whether presenting local taxonomy structure together with instructions for interpreting it improves that task for large language models. The central comparison holds the entity, description, and prompting style fixed while changing both structural evidence and graph-specific instructions. Three styles are examined: standard prompting, chain-of-thought prompting, and contrastive chain-of-thought prompting.


The study uses the released 1,000-entity WiKC evaluation sample and the refined Wikidata taxonomy introduced by Peng, Bonald, and Alam (2024). Nine conditions are evaluated with Gemini 3.8 Flash and Llama 3.1 8B Instruct; Llama is run with seeds 0, 42, and 123. The supplied archives contain 9,000 completed Gemini responses and 8,991, 8,971, and 8,986 completed main Llama responses, respectively. A further child-removal experiment at Llama seed 0 completes 1,352 of 1,353 planned responses for 451 eligible entities. Within-model contrasts use matched successful responses. Context is selected around an annotated gold class, whose label is omitted; the experiment therefore measures the value of oracle structural evidence rather than autonomous retrieval.


Graph-aware prompting improves mean entity F1 in every model-by-style comparison, with paired gains from 0.0451 to 0.1348. Parent-only context outperforms sibling-only context on aggregate F1. The three Llama seeds permit a direct check of sampling variation, while the matched child intervention separates child-context effects from differences in cohort composition. The results are interpreted together with output length, ranking metrics, and condition-specific completion rates.


The empty-child control gives small positive point estimates (about +0.005 to +0.024); Gemini’s roughly +0.02 contrasts with its +0.0625 standard graph gain. This suggests a template contribution, but different cohorts prevent attributing a fixed share of the full gain to instructions.


These findings support graph-aware prompting as a useful source of evidence under controlled context selection. The contribution is a reproducible prompting and evaluation framework, together with an empirical account of where local structure helps and where it misleads. Future work should evaluate retrieved context, independently annotated gold sets, and controlled structural noise.


{\small Keywords: entity typing, knowledge graphs, taxonomy, prompting, WiKC, evaluation\par}

